% !PRINTED 354
\weberpage{354}

Since $a$ is to be divisible by $\mathfrak{a}$, we set
\[ a = \mathfrak{a}\mathfrak{m}' = \mathfrak{a}'\mathfrak{m} \]
and obtain from (1):
\weq{2}{\mathfrak{a} = m\mathfrak{m}.}
Thus $m$ divides the divisor of the ideal $\mathfrak{a}$, and with this Theorem 3. is proved, since, if $\mathfrak{a}$ is primary, its divisor must be $= 1$.

4. \gesperrt{Theorem. If $\mathfrak{m}$ is a primary ideal of the field $\Omega$, then every integer in $\Omega$ is congruent to a rational number with respect to the modulus $\mathfrak{m}$.}

For by Bd.\ II, \S\ 165, $N(\mathfrak{m})$ is in general the number of the numbers in $\Omega$ that are incongruent with respect to $\mathfrak{m}$. If now $N(\mathfrak{m}) = a$ is at the same time the least \gesperrt{natural} number divisible by $\mathfrak{m}$, then the $a$ numbers
\[ 0, 1, 2, \ldots, a - 1 \]
are all incongruent, and among them all classes of numbers modulo $\mathfrak{m}$ are represented.

Let now $[Q]$ be an order with the conductor $Q$ and the discriminant $D$, and $(1, \theta)$ the basis of this order. Further let $\mathfrak{m}$ be a \gesperrt{primary} ideal coprime to $Q$. By Theorem 4. there is then a rational integer $t$ which satisfies the condition
\[ \theta + t \equiv 0 \pmod{\mathfrak{m}} \]
We set
\weq{3}{\begin{aligned}
t &= \frac{b}{2}, && \text{if } D \equiv 0 \\
t &= \frac{1+b}{2}, && \text{„} \quad D \equiv 1
\end{aligned}
\qquad \text{(mod 4)}}
and obtain from \S\ 96, (5):
\weq{4}{\theta + t = \frac{b + \sqrt{D}}{2} \equiv 0 \pmod{\mathfrak{m}},}
wherein $b$ is a rational integer which by (3) satisfies the condition
\[ b \equiv D \pmod{2} \]
Moreover let
\weq{5}{N(\mathfrak{m}) = a,}
so that $a$ is the least natural number divisible by $\mathfrak{m}$.

By the congruence (4) the number $b$ is determined only with respect to the modulus $2a$. For if $b, b'$ are two numbers which satisfy the condition (4), then $\frac{1}{2}(b' - b)$ is divisible by $a$.

Since further
\[ N\left(\frac{b + \sqrt{D}}{2}\right) = \frac{b^2 - D}{4} \]
%CONT>>
